\chapter{Transition Metals: Color, Magnetism, and Catalysis}

\begin{kaichang}
Find four objects: a rusty nail, reddish copper stripped from discarded wire, a shiny stainless-steel spoon, and a small black refrigerator magnet.

Test your eye. The nail and spoon mainly belong to the same periodic-table neighborhood from Chapter 10. Why is one rusty and the other bright after ten years? Why is copper reddish brown while most metals are silver-white? And why, among over a hundred elements, are iron, cobalt, and nickel the familiar metals strongly attracted by magnets? These brothers are neighbors in the table. Coincidence?

One more puzzle: a laboratory blue crystal called blue vitriol dissolves into a beautiful sapphire-blue solution. Its ingredients are merely copper, sulfur, oxygen, and water. Most sulfates are colorless, so who dyed this solution blue?

All four answers lie in the central district we have not yet explored carefully. This chapter knocks on its doors one by one.
\end{kaichang}

\section{Three Specialties of the Central District}

Chapter 10 called transition metals the central low-rise neighborhood: iron, copper, zinc, titanium, chromium, manganese, cobalt, and nickel, with silver, gold, platinum, and mercury below. It offers three specialties unavailable from the districts on either side.

First, \textbf{color}. Main-group substances are plain: salt, sugar, quartz, and soda ash are nearly all white or colorless. Transition metals form a pigment shop: blue copper-ion solutions, deep purple potassium permanganate antiseptic, green chromium-ion solutions, and reddish brown rust. Jewelers know this well. Ruby and sapphire share the same colorless aluminum-oxide framework, corundum, but ruby contains a few parts per ten thousand chromium atoms, sapphire iron and titanium. Tiny transition-metal impurities turn colorless stone into national treasures.

Second, \textbf{magnetism}. Bring elemental samples to a magnet and most do nothing, while iron, cobalt, and nickel cling strongly and can become magnets themselves. Refrigerator magnets, headphone speakers, and electric-car motors depend on these brothers.

Third, \textbf{catalysis}. Catalysts accelerate reactions without appearing among final products. Fertilizer ammonia synthesis uses iron; expensive automotive three-way catalytic converters use platinum, palladium, and rhodium; black manganese dioxide added to hydrogen peroxide features manganese. Transition metals and their compounds occupy much of industrial and agricultural catalyst lists.

Three specialties, three clues: we will see their common root.

\section{Ten Extra Boxes: Enter the d Electrons}

Return to the table's shape. Chapter 10 linked the first row's 2 and next two rows' 8 boxes to electron floors. Row four suddenly widens from 8 to 18. Its 10 extra boxes, scandium through zinc, are the first transition-metal row.

Those boxes represent new d rooms: five per shell, two electrons each, exactly 10 places. Chapter 10's challenge mentioned the crucial twist: fourth-shell s rooms are cheaper than third-shell d rooms, so row-four elements fill 4s first, then return to 3d. Added electrons from scandium to zinc occupy these comeback rooms.

Why does this matter? It creates a special situation: \textbf{3d and 4s electrons have almost equal energies}. In main-group atoms, widely separated energies hold inner electrons firmly out of chemistry, leaving only the outer one or two active. Here, 3d and 4s are so close that chemical energies can recruit them in turn.

Consequently, \textbf{one element can surrender different electron counts and exhibit several oxidation states}. Iron loses two for \chem{Fe^{2+}} or another for \chem{Fe^{3+}}; copper gives \chem{Cu^{+}} and \chem{Cu^{2+}}. Manganese is dramatic, offering seven or eight valences from $+2$ through $+7$, with permanganate \chem{MnO_4^{-}} at $+7$. Main-group sodium always gives one electron, one lifelong valence, while manganese plays many roles. Variable states introduce varied colors and catalytic power.

\section{Where Color Comes From: The Light Absorbed}

First ask what color means. White sunlight or lamplight mixes colors. An object's appearance depends on which it \textbf{absorbs}. Leaves absorb red and blue-violet, reflecting green to your eyes. Asking why copper sulfate is blue is asking who absorbs orange-yellow.

d electrons answer. Their five 3d rooms are not all equal in energy. An electron can jump from lower to higher if incoming light exactly matches the gap, absorbing that light. Crucially, \textbf{d-room energy gaps fall within visible-light energies}. Just right, neither more nor less. Transition-metal ions absorb part of visible light and return the remaining colors.

Compare sodium and calcium ions. They too have electron rooms, but empty rooms lie much farther away, with ultraviolet energy gaps. Human eyes do not see ultraviolet, so absorbing it leaves solutions colorless. Transparent saltwater absorbs light, just light you cannot see.

An even better detail: d-room gaps are not fixed but \textbf{altered by surrounding objects}. This leads to coordination, the chapter's central concept.

\section{Coordination: Tenants Surround the Metal Ion}

\chem{Cu^{2+}} never stands alone in water. Water oxygens have spare electron pairs attracted to positive copper. Four to six molecules gather and lend their lone pairs, forming a stable surrounding group. The whole is a \textbf{coordination compound}, or complex: a central metal ion surrounded by electron-pair-donating molecules or ions, \textbf{ligands}. Strictly, copper sulfate's blue belongs to copper plus six water ligands, not \chem{Cu^{2+}} alone.

The evidence is easy to demonstrate. Anhydrous copper sulfate, stripped of water ligands, is gray-white powder. Water immediately surrounds it and restores blue. Change ligands and color changes: concentrated ammonia added to blue copper solution replaces water and produces much deeper blue-violet. This releases irritating gas; watch a teacher or video, not a home experiment. One copper ion, different tenants, different d-energy splitting, different absorbed light, different color.

This explains many phenomena. Chromium in ruby's aluminum-oxide lattice is surrounded by oxygens and gives red. Chromium in emerald's beryl lattice has a different surrounding arrangement and gives green. Ligands essentially tune chromium's gemstone colors.

\section{Magnetism: Billions of Tiny Needles Lined Up}

What about magnetism? d electrons remain the source, but now the keyword is \textbf{spin}.

Chapter 9 introduced spin, pictured as an electron behaving like an extraordinarily tiny magnetic needle. This aids thought; the electron is not literally spinning. Paired electrons in one room point their needles oppositely and cancel. Transition-metal d rooms are often partly filled: \chem{Fe^{3+}} has five d electrons, each singly occupying a room under Hund's rule, all five needles aligned without cancellation. Such an ion is a miniature magnet.

Miniature magnets alone are insufficient. Ordinary iron on a table does not attract nails because it contains countless \textbf{magnetic domains}. Within each, billions of needles align, but domains point randomly and cancel. Iron, cobalt, and nickel have interactions making neighboring-needle \textbf{alignment especially inexpensive}. An external magnet makes domains turn together as if commanded, magnetizing the whole piece strongly. Copper, aluminum, and zinc either cancel paired d electrons or lack cooperative alignment, so miss out on magnetism. The brothers' neighboring positions are no coincidence: similar 3d occupancies satisfy the demanding combination of many needles and easy alignment.

\section{Rust, Stainless Steel, and a Slow Electrochemical Fire}

Now the first opening puzzle: why does iron rust while stainless steel does not?

A nail in dry air can last years; in boiled, sealed water, it also lasts. When both water and oxygen arrive, red-brown rust appears. Iron transfers electrons to oxygen, first becoming \chem{Fe^{2+}}, with electrons moving through metal and water film to oxygen. Subsequent steps form hydrated iron oxide, whose main composition can be written \chem{Fe_2O_3}:
\[\chem{4Fe} + \chem{3O_2} \rightarrow \chem{2Fe_2O_3}\]
Real rust contains variable water and is porous like a soaked sponge. That porosity is fatal: it cannot block further water and oxygen, so corrosion eats inward layer after layer until the nail fails.

Stainless steel remains mainly iron but contains roughly $12\%$ or more chromium. Chromium binds oxygen more strongly and forms chromium oxide first. Unlike rust, this film is \textbf{dense, strong, and transparent}. Only a few nanometers thick, it excludes water and air. Better still, it heals: scratching exposes chromium that immediately oxidizes and repairs the film. Invisible, it guards constantly. More precisely, stainless steel rusts so quickly and so well that it makes perfect armor and stops.

Corrosion defenses extend this logic. Paint and oil physically exclude oxygen and water. Zinc plating adds strategy: zinc gives electrons more readily than iron, so even damaged plating sacrifices zinc first, sacrificial-anode protection. Ships and bridge piles carry welded zinc or magnesium blocks, replaced periodically.

\begin{shiyan}
\textbf{Compare Nail Rusting} (Safe at home throughout.)

Materials: four clean nails, four transparent cups, tap water, salt, cooking oil, and paper towels.

Procedure: (1) Put a nail in open cup one with dry air only and dry paper beneath it. (2) Half-submerge a nail in tap water in cup two, leaving half exposed. (3) Do the same in cup three with two spoonfuls of dissolved salt. (4) Fully cover a nail with tap water in cup four and seal the surface with oil.

What to observe: Check and record daily for three to seven days. Which rusts fastest, and where do spots begin: underwater or near the waterline? Cup two's wet--dry boundary rusts visibly, saltwater cup three faster, while oil-covered four and driest one barely rust.

Draw the conclusion yourself: rust requires water and oxygen \textbf{together}, and salt intensifies this slow fire. You now know why seawater ships rust readily.
\end{shiyan}

\section{Catalysts: Building a Tunnel Through the Mountain}

Catalysis is our third specialty. Chapter 27's energy language separates whether and how far a reaction proceeds, thermodynamics, from how quickly, kinetics. The key obstacle is \textbf{activation energy}, the mountain reactants must cross before becoming products. Higher mountains permit fewer crossings and slower reactions.

A catalyst does one thing: \textbf{provides a route with a lower mountain}. It does not push reactants, supply energy, or make an impossible reaction possible; tunneling cannot help a fundamentally uncrossable mountain. It simply makes a possible but impractically slow reaction useful. Starting and ending points stay fixed, so released heat and final equilibrium are unchanged.

Why are transition metals good tunnel builders? Again, d electrons and variable states. Consider surface-catalyzed ammonia synthesis,
\[\chem{N_2} + \chem{3H_2} \rightleftharpoons \chem{2NH_3}\]
whose main obstacle is nitrogen's triple bond. Iron-surface d electrons first capture, or adsorb, \chem{N_2}. Multiple surface hands hold nitrogen atoms and weaken and split the bond. Hydrogen likewise splits, then atoms meet sequentially on the workbench to assemble ammonia and depart. Breaking, holding, arranging, and matchmaking open a previously unavailable low-energy route. Iron returns to its original surface state for the next guests. Platinum, palladium, and rhodium do similar work in automotive converters, capturing carbon monoxide, nitrogen oxides, and unburned hydrocarbons and turning them into carbon dioxide, nitrogen, and water. A small exhaust canister employs some of the world's costliest grams of metal.

\begin{qiaoliang}
How much faster is faster? Activation-energy effects are exponential: under similar conditions, rate is roughly proportional to $e^{-E_a/(RT)}$, with activation energy $E_a$, temperature $T$, and constant $R$. Near room temperature, $RT$ is about \ud{2.5}{kJ\,mol^{-1}}. A modest catalytic reduction of \ud{20}{kJ\,mol^{-1}} increases rate by about
\[e^{20/2.5} = e^{8} \approx 3000\ \text{times}\]
A \ud{40}{kJ\,mol^{-1}} reduction gives $e^{16}$, about ten million times. Small barrier decreases \textbf{multiply} rather than add to rates. Thus new catalysts often cost less than heating and compression in chemical plants. Life's enzymes, protein catalysts starring in Volume II, play this exponential game to perfection.
\end{qiaoliang}

\begin{wuqu}
``Catalysts participate, so they are gradually consumed; declining performance means they are used up.'' Participation is true, consumption false. They capture molecules, break bonds, and arrange partners, but \textbf{return to their original state} each cycle and are theoretically not consumed. Industrial decline usually means poisoning or wear: impurities such as sulfur occupy the workbench, or heat sinters surfaces and reduces area. The workbench is dirty or damaged, not used up. A laboratory spoonful of manganese dioxide catalyzing peroxide can be filtered, dried, and weighed afterward with no mass lost, ready for reuse.
\end{wuqu}

\section{How Scientists Know: A Century's Debate About Color}

\begin{zhengju}
Late nineteenth-century chemists had puzzling compounds. Ammonia added to cobalt salts yielded \chem{CoCl_3 \cdot 6NH_3}, \chem{CoCl_3 \cdot 5NH_3}, and \chem{CoCl_3 \cdot 4NH_3}, equal cobalt with declining ammonia but yellow, purple, and green colors. Chain-bond theories could not draw where ammonia attached beside cobalt chloride.

In 1893, Swiss chemist Werner proposed an initially heretical model: cobalt at the center, ammonia and chloride \textbf{surrounding it directly} in an inner sphere, extra chloride outside balancing charge. A falsifiable prediction followed: inner chloride would remain cobalt-bound in water, while outer chloride became free. Dissolve samples, add silver nitrate, and measure immediate silver chloride precipitation. The 6-ammonia compound yielded 3 chloride portions, the 5-ammonia compound 2, the 4-ammonia compound 1, exactly as predicted. Alternatives, such as ammonia clusters hanging from chains, fell one by one to conductivity and precipitation measurements. Stronger evidence followed: Werner predicted non-superimposable mirror-image complexes like hands and separated them himself in 1911, nearly closing every escape route. He received the 1913 Nobel Prize in Chemistry.

Remember two points. Coordination theory prevailed through twenty years of quantitative predictions and experiments eliminating rivals, not one genius experiment. And Werner knew nothing of d electrons: electrons were newly discovered and quantum mechanics unborn. Colors and precipitates let him infer spatial atomic arrangements. Finding patterns before mechanisms is normal science, no embarrassment.
\end{zhengju}

\section{Life Used Transition Metals First and Best}

Thinking coordination and catalysis belong only to factories and jewelers underestimates nature. Life recruited all three specialties over three billion years ago.

Start with blood. Hemoglobin gives it red color, and at each molecule's center sits a complex: an iron ion embraced from four directions by the ring ligand heme, with a remaining position for oxygen. \chem{Fe^{2+}} binding \chem{O_2} must be \textbf{just right}. Ligands tune it neither too strong nor weak, capturing abundant lung oxygen and releasing it in oxygen-poor tissues. Carbon monoxide's danger is occupying the same site over two hundred times more tightly, refusing to leave and filling hemoglobin's vehicles. Why iron? Its d electrons and variable states suit reversible oxygen capture and release, a specialty main-group metals cannot provide.

Now cobalt. Vitamin \chem{B_{12}}, the most structurally complex known vitamin, has a central cobalt ion in a ring ligand. Humans cannot make it and rely on intestinal bacteria and food; deficiency disrupts blood formation and nerves. Plants also benefit: legume-root nitrogenase uses molybdenum and iron centers to fix nitrogen at ordinary temperature and pressure, where factories need extremes. Its exact mechanism still has unresolved details; science has not finished this exam.

Estimate the scale. An adult contains about \ud{4}{g} of iron, mostly in hemoglobin. With molar mass \ud{56}{g\,mol^{-1}}, $4/56 \approx 0.07$ mol, multiplied by \ud{6.02\times 10^{23}}{} atoms per mole gives
\[0.07 \times 6.02\times 10^{23} \approx 4\times 10^{22}\ \text{iron atoms}\]
Four hundred quintillion iron atoms take turns capturing and releasing oxygen in your blood now. Every breath depends on transition-metal coordination working properly.

\section{The Model's Boundaries}

As usual, useful explanations must be followed by their failure modes.

First, d-room transitions explain many colors but not all. Permanganate \chem{MnO_4^{-}} has manganese at $+7$ and empty d rooms, no d electrons to jump. Its deep purple comes from ligand-to-metal charge transfer, electrons moving from oxygen toward manganese, which absorbs more strongly. All color from d transitions is a useful first approximation, not the complete theory.

Second, electron rooms and magnetic needles are models throughout. They explain trends and exceptions well, but electrons are not literal needles, and domain alignment involves subtler quantum exchange interactions. For exact alloy magnetism, this picture points the way without finishing the calculation.

Third, surface catalysis has been greatly simplified. Real surfaces are irregular and often only a few active sites matter. Impurities, temperature, and particle size can overturn expectations. Iron catalyzes ammonia is true; why one piece works better than another remains a research frontier.

\section{Back to the Four Objects}

Check the answers. Rusty nail and stainless spoon are both mainly iron, differing in chromium. Chromium quickly forms a dense, transparent oxide excluding water and oxygen; ordinary rust is porous and water-loving, inviting trouble inward.

Copper wire is reddish because its 3d rooms are nearly but not completely full. d transitions absorb blue-violet and reflect red-orange. Silver's corresponding gap lies in ultraviolet, leaving visible light intact and giving silver-white. Metallic color is another d-electron account.

The magnet uses iron, cobalt, or nickel's partly filled 3d rooms, many unpaired electrons, and easy cooperative alignment. Both conditions are necessary, so the magnetic brothers' neighboring positions are no coincidence.

Blue vitriol has copper surrounded by water ligands that tune d gaps to absorb orange-yellow, leaving blue-violet for your eyes. Remove water by heating and color vanishes; substitute ammonia and it deepens. The blue is a joint signature of copper and six water tenants.

\section{The District Still Has More to Explore}

We have explained transition-metal colors, magnetism, and catalysis while raising two larger questions. Where did iron, cobalt, nickel, and copper originate? Crustal iron is only a small share compared with the core, and your 4 grams, motor cobalt, and phone rare earths traveled astonishing journeys. Chapter 16 follows stellar furnaces to blood vessels. And why do enzymes at body temperature work faster and better while iron-catalyzed ammonia needs high temperature and pressure? How proteins amplify transition-metal abilities by orders of magnitude awaits Volume II, after carbon's construction of life's large molecules.

\begin{tiaozhan}
Why do d-transition gaps lie in visible energies rather than deep ultraviolet or infrared? Roughly, d electrons are not much farther from nuclei than s and p electrons, and ligand-pair repulsion produces splitting around \ud{100}{} to \ud{400}{kJ\,mol^{-1}}. Convert photon energy $E = hc/\lambda$ to molar terms: shorter $\lambda$ means higher energy. Visible limits $\lambda \approx 400$ nm and $700$ nm give about \ud{300}{} to \ud{170}{kJ\,mol^{-1}}, nearly overlapping ligand-field splitting. This is no designed coincidence, but electromagnetic strength at atomic scales. Tenfold larger splitting would make every complex colorless; tenfold smaller would make them all ink-black. Why five d rooms split into three lower and two higher under octahedral ligands requires drawings of orbital directions and ligand positions, the opening of university ligand-field theory. Skipping this does not affect the main story.
\end{tiaozhan}

\section*{Try It Yourself}

\begin{sikaoti}
\item Draw coordinated \chem{Fe^{2+}} in water, six droplet-shaped waters around central iron with oxygen ends inward. Identify the inner sphere and explain why this whole has color while saltwater does not.
\item Iron rusts in damp air but survives millennia in dry deserts. Draw water, oxygen, and iron's roles in a material-flow map, emphasizing particles required together. Explain severe ship corrosion near waterlines.
\item Draw energy versus reaction progress with reactants, products, and a mountain. Add a dashed catalyzed route. Identify activation energy and what catalysis changes or preserves.
\item Explain ruby and sapphire, the same colorless aluminum-oxide framework with trace impurities, using coordination and d electrons. Why do different transition metals give different colors, and why can tiny concentrations be conspicuous?
\item Someone says stainless steel never rusts because it contains no rustable components. Identify both mistakes and rewrite the claim using chromium and oxide films.
\item Inspect kitchen pots, spoons, taps, and cans. List transition-metal uses and guess the property involved: color, magnetism, catalysis, dense oxide film, thermal conduction, or another. Record any uncertain case and see whether later chapters answer it.
\end{sikaoti}
